\section{Методология численного эксперимента}
\subsection{Цель эксперимента}
Эксперимент направлен на проверку гипотезы о том, что нелинейные модели
превосходят линейные в задачах семантического сходства.

\begin{table}[h]
  \centering
  \caption{Спецификация моделей}
  \begin{tabularx}{\linewidth}{lX}
    \toprule
    \textbf{Компонент} & \textbf{Параметры} \\
    \midrule
    \multicolumn{2}{l}{\textit{LSA-конвейер}} \\
    \quad TF-IDF      & \texttt{sublinear\_tf=True},
    \texttt{min\_df=2}, \texttt{max\_df=0.95} \\
    \quad TruncatedSVD & $k \in \{50, 100, 200, 300, 500\}$,
    \texttt{random\_state=42} \\
    \quad Нормализация & L2 \\
    \midrule
    \multicolumn{2}{l}{\textit{all-MiniLM-L6-v2}} \\
    \quad Архитектура  & 6 слоёв, $d = 384$ \\
    \quad Агрегация    & Mean pooling с маской внимания \\
    \quad Нормализация & L2 \\
    \bottomrule
  \end{tabularx}
\end{table}

\subsection{Датасеты}

\textbf{STS Benchmark (STSb)} ~--- стандартный корпус из 8 628 пар
предложений с
оценками семантической близости от 0 до 5. Используется тестовая часть (1 379
пар).
\textbf{20 Newsgroups} ~--- коллекция новостных постов из 20 тематических
категорий. Выбраны 10 категорий:
\begin{itemize}
  \item \texttt{sci.med}, \texttt{sci.electronics}, \texttt{sci.space}
  \item \texttt{rec.sport.hockey}, \texttt{rec.sport.baseball}
  \item \texttt{talk.politics.guns}, \texttt{talk.politics.mideast}
  \item \texttt{comp.graphics}, \texttt{comp.os.ms-windows.misc}
  \item \texttt{alt.atheism}
\end{itemize}
В выборке имеются как тематически близкие (\texttt{sci.med,
sci.space}), так и контрастные
(\texttt{alt.atheism, rec.sport.hockey}) категории.

\textbf{MRPC} ~--- корпус из 5 801 пар предложений с бинарными
метками парафраз
(3 900 ~--- парафразы, 1 901 ~--- не парафразы).

\subsection{Метрики}
\subsubsection{Интринсивные метрики}
Данные метрики характеризуют геометрию пространства эмбеддингов и не зависят
напрямую от внешней задачи.

\paragraph{1. Participation Ratio (PR)}

\begin{displaymath}
  \text{PR}(\mathbf{E}) = \frac{\left(\sum_{i}
  \lambda_i\right)^2}{\sum_{i} \lambda_i^2}
\end{displaymath}
где \(\lambda_i\) ~--- собственные значения матрицы ковариаций
\(\mathbf{C} = \frac{1}{n}\mathbf{E}^\top \mathbf{E}\) (после
вычитания среднего), \(\mathbf{E} \in \mathbb{R}^{n \times d}\) ~---
матрица эмбеддингов. Высокий PR означает, что дисперсия распределена
по многим направлениям, а не сосредоточена в нескольких главных компонентах.

\paragraph{2. Effective Rank (ER)}

\begin{displaymath}
  \text{ER}(\mathbf{E}) = \exp\!\left(-\sum_{i} p_i \log p_i\right),
  \quad p_i = \frac{\sigma_i}{\sum_j \sigma_j}
\end{displaymath}

где \(\sigma_i\) — сингулярные числа \(\mathbf{E}\), \(p_i\) —
нормированные веса, формирующие дискретное распределение. ER — это
экспонента энтропии Шеннона спектра сингулярных чисел; максимально
при равномерном распределении. Высокий ER свидетельствует о том, что
информация распределена по многим сингулярным направлениям.

\paragraph{3. Isotropy Score (IS)}

\begin{displaymath}
  \text{IS}(\mathbf{E}) = \frac{1}{|I|} \sum_{\omega \in I}
  \frac{\min_{\mathbf{v}} \exp\!\left(f(\omega)^\top
  \mathbf{v}\right)}{\max_{\mathbf{v}} \exp\!\left(f(\omega)^\top
  \mathbf{v}\right)}
\end{displaymath}

IS = 1 при идеальной изотропии (равномерном распределении на
гиперсфере), \(\text{IS} \to
0\) при концентрации в низкоразмерном подпространстве. В коде используется
упрощенная версия: \(\text{IS} \approx \lambda_{min} / \lambda_{max}\).

\paragraph{4. Hubness Score (HS)}

\begin{displaymath}
  \text{HS} = \frac{S_{N_k}}{\mu_{N_k}}
\end{displaymath}

где \(S_{N_k}\) ~--- стандартное отклонение, \(\mu_{N_k}\) ~--- среднее число
вхождений точки в \(k\text{-NN-списки}\). Высокий HS \(> 1\) указывает на
концентрацию \enquote{хабов}: точек в пространстве, которые
оказываются ближайшим соседом для непропорционально большого числа
других точек.

\subsubsection{Экстринсивные метрики}

\paragraph{1. Оценка семантического сходства (задача STS).}

Корреляция Спирмена (\(\rho\)) оценивает монотонность
связи между предсказаниями модели и человеческими оценками, не требуя
линейности. Формула вычисляется через разность рангов \(d_i\) для
каждой из \(n\) пар:

\begin{displaymath}
  \rho = 1 - \frac{6\sum_{i=1}^{n} d_i^2}{n(n^2 - 1)}
\end{displaymath}

Значение \(\rho\) близкое к 1 свидетельствует о том, что модель
ранжирует пары по сходству так же, как и человек.

\paragraph{2. Классификация.}

Основной метрикой является Accuracy (доля
верных предсказаний). Для учета возможного дисбаланса используется
F1-score (гармоническое среднее точности P и полноты R):

\begin{displaymath}
  \text{F1} = 2 \cdot \frac{P \cdot R }{P + R}
\end{displaymath}

Используется macro-усреднение: F1 вычисляется независимо для каждого
из \(C\) классов, затем берётся среднее арифметическое:

\begin{displaymath}
  F1_{\text{macro}} = \frac{1}{C}\sum_{c=1}^{C} F1_c
\end{displaymath}

что позволяет оценить качество модели на каждой тематике равновесно.

\paragraph{3. Кластеризация.}

\begin{enumerate}
  \item \textbf{Silhouette Score (коэффициент силуэта):} Измеряет,
    насколько объект \(i\) похож на свой кластер по сравнению с другими
    кластерами. Для каждой точки рассчитывается среднее расстояние до
    точек своего кластера (\(a\)) и до точек ближайшего соседнего
    кластера (\(b\)):
    \[
      s(i) = \frac{b(i) - a(i)}{\max\{a(i),\, b(i)\}}
    \]
    Среднее значение по всем документам в интервале \([-1, 1]\) показывает
    качество разбиения: значения выше 0.5 указывают на наличие выраженной
    структуры.

  \item \textbf{Adjusted Rand Index (ARI):} Оценивает сходство
    полученных кластеров с истинными тематическими метками, корректируя
    результат на вероятность случайного угадывания. Значение 1
    соответствует идеальному совпадению с категориями датасета.
\end{enumerate}

\paragraph{Проверка парафраз (датасет MRPC)}
В задачах определения парафраз классы часто не сбалансированы (парафразы
встречаются реже), поэтому F1-score является более надежным
индикатором, чем точность. Высокий балл подтверждает, что модель
способна отличать тонкие семантические различия от простого пересечения слов.
